\documentclass{article}
\usepackage[utf8]{inputenc}
\usepackage{amsmath, amssymb, amsthm, mismath, ulem, gensymb, mathtools, nccmath, hyperref, tikz, pgfplots}
\usepackage{gnuplottex}
\usepackage{lipsum}
\usepackage{xcolor}
\usepackage{etoolbox}
\usepackage{hyperref}
\usepgfplotslibrary{fillbetween}
\pgfplotsset{compat=1.18}

\definecolor{desgreen}{HTML}{348443}
\DeclareMathOperator{\erfc}{\mathrm{erfc}}

\newcommand{\dx}{\,\mathrm{d}x}
\newcommand{\dt}{\,\mathrm{d}t}
\newcommand{\ddx}{\,\frac{\mathrm{d}}{\mathrm{d}x}}
\newcommand{\dydx}{\, \frac{\mathrm{d}y}{\mathrm{d}x}}
\newcommand{\du}{\, \mathrm{d}u}
\newcommand{\dv}{\, \mathrm{d}v}
\newcommand{\dw}{\, \mathrm{d}w}
\newcommand{\dy}{\, \mathrm{d}y}
\newcommand{\dr}{\, \mathrm{d}r}
\newcommand{\dtheta}{\, \mathrm{d}\theta}

\let\oldint\int
\renewcommand{\int}{\oldint\limits}

\delimiterfactor 1200




\newcounter{problem}
\newcommand{\problem}[1]{
	\refstepcounter{problem}
	\noindent\textbf{\theproblem}\quad #1
}

\newtoggle{solutions}
\togglefalse{solutions}

\newtoggle{studentpagebreaks}
\newtoggle{studentvspace}

\newtoggle{instructorpagebreaks}
\newtoggle{instructorvspace}

\iftoggle{solutions}{
	\togglefalse{studentpagebreaks}
	\togglefalse{studentvspace}
	\toggletrue{instructorpagebreaks}
	\toggletrue{instructorvspace}
}{
	\toggletrue{studentpagebreaks}
	\toggletrue{studentvspace}
	\togglefalse{instructorpagebreaks}
	\togglefalse{instructorvspace}
}

\newcommand{\spb}{
	\iftoggle{studentpagebreaks}{\newpage}{}
}
\newcommand{\svs}{
	\iftoggle{studentvspace}{\vspace*{0.5\textheight}}{}
}
\newcommand{\ipb}{
	\iftoggle{instructorpagebreaks}{\newpage}{}
}

\title{Gamma Function}
\author{By SMC Integration Club}
\date{\today}

\hypersetup{
	pdftitle={Gamma Function - SMC Integration Club},
	pdfauthor={Thomas Yamada, Nicholas Quiles},
	pdfsubject={Integration Club Gamma Function Worksheet},
	pdfkeywords={calculus, integration, gamma function}
}

\begin{document}
	\maketitle
	
	{\centering
		What is the Gamma Function? \\
		[12pt]
		The Gamma Function is defined as: 
		\begin{equation}
			\Gamma(z) = \int_{u=0}^{\infty} u^{z-1} e^{-u} \dt
		\end{equation}
		\text{Graph of }$\Gamma(z)$:\par
	}
	\vspace{-10pt}
	\begin{figure}[!hbtp]
		\centering
		\begin{tikzpicture}
			\begin{axis}[
				axis lines = center,
				xmin = -5,
				xmax = 5,
				ymin = -10,
				ymax = 10,
				minor tick num=4,
				minor tick style={draw=none},
				grid = both,
				width = \textwidth,
				height = 0.5\textwidth,
				axis line style={draw=black, thick},
				major tick length=5pt,
				ticklabel style={font=\bfseries},
				tick style={thick},
				xlabel = {$x$},
				ylabel = {$y$},
				unbounded coords=jump,
				clip=true,
				]
				\addplot[very thin, dashed, desgreen] coordinates {(0, 10) (0, -10)};
				\addplot[very thin, dashed, desgreen] coordinates {(-1, 10) (-1, -10)};
				\addplot[very thin, dashed, desgreen] coordinates {(-2, 10) (-2, -10)};
				\addplot[very thin, dashed, desgreen] coordinates {(-3, 10) (-3, -10)};
				\addplot[very thin, dashed, desgreen] coordinates {(-4, 10) (-4, -10)};
				\addplot[line width=1.2pt, red, samples=1600] gnuplot {
					(x < 0.2 && abs(sin(pi*x)) < 0.02) ? NaN : gamma(x)
					};
			\end{axis}
		\end{tikzpicture}
		\caption{Graph of the Gamma Function $\Gamma(z)$.}
	\end{figure}
	You can see, by the graph, that the function is neither odd or even, because it's asymmetric everywhere. There is no linearity either.\\
	[5pt]
	
	There is two relations though. It's core identity is the recurrence relation, the Gamma's equivalent of a "linearity rule":
	
	\begin{equation}
		\Gamma(z+1) = z\Gamma(z)
	\end{equation}
	
	Then, you have Euler's Reflection Formula, a sort of complementary behavior:
	\begin{equation}
		\Gamma(z)\Gamma(1-z) = \frac{\pi}{\sin(\pi z)}
	\end{equation}
	
	So, what is this Gamma Function really? Remember the factorial? Well,
	\begin{equation}
		\Gamma(z+1) = z!
	\end{equation}
	
	It's partially that simple, but also not exactly. If you remember using the factorial, you've only ever used integers, because that's easy to compute with the basic function. i.e. $6! = 6\cdot5\cdot4\cdot3\cdot2\cdot1 = 720$ \\
	
	But, how would you do that with $\frac{1}{2}!$? What about $-\frac{1}{2}!$? That's where the Gamma Function steps in. It allows you to computer the "factorial" of non integers. It does have some restrictions though:
	
	\begin{equation}
		\mathrm{Dom}(\Gamma) = \mathbb{R} \backslash \left\lbrace -n : n \in \mathbb{Z}_{\geq 0} \right\rbrace
	\end{equation}
	
	\noindent Let's try $\Gamma(1)$:
	\begin{equation}
		\begin{aligned}
		\Gamma(1) &= \int_{u=0}^{\infty}u^{0}e^{-u}\du = \int_{u=0}^{\infty}e^{-u}\du \\
		\Gamma(1) &= \left[-e^{-u}\right]_{u=0}^{\infty} = 0 - (- 1) \\
		\Gamma(1) &= 1 = 0!
		\end{aligned}
	\end{equation}
	\noindent For $\Gamma(2)$, you'll need to use a simple IBP: \\[5pt]
	\begin{equation}
		\begin{aligned}
			\Gamma(2) &= \int_{u=0}^{\infty}u^{1}e^{-u}\du = \int_{u=0}^{\infty}te^{-u}\du \\
			\Gamma(2) &= -ue^{-u} + \int_{u=0}^{\infty}e^{-u}\du = 0 + 1 \\
			\Gamma(2) &= 1 = 1!
		\end{aligned}
	\end{equation}
	\noindent You can probably see where this is going. $\Gamma(3) = 2 = 2!$ \\[10pt]
	So, what about the recurrence relation? How can we use that? Well, keep in mind that computing $\Gamma(0)$ does not exist, as you can see from the graph above. So, if we were to apply $\Gamma(z+1) = z\Gamma(z)$ to get, for example, a negative integer, you'd need to re-use the relation again and again until you reached a computable number, but you'd always hit $0$ before anything actually computable, which is why every negative integer less than $0$ does not exist. \\
	
	Then, what about negative fractionals? Let's use some negative fractional that is perfectly divisible by $\frac{1}{2}$, such as $-\frac{3}{2}$. You'll need to use the recurrence relation in order to get to that. The first directly computable value for that is $\Gamma(\frac{1}{2})$, so we need to solve that first. Luckily, it gives a nice number to work with.
	
	\newpage
	Let's look into how we'd find $\Gamma\left(\frac{1}{2}\right)$: \\[10pt]
	\[
	\displaystyle{\Gamma\left(\frac{1}{2}\right) = \int_{u=0}^{\infty} e^{-u}u^{-\frac{1}{2}}\du}
	\] \\ [5pt]
	\noindent Let's rewrite it so it's easier to read.
	\[\Gamma\left(\frac{1}{2}\right) = \int_{u=0}^{\infty} e^{-u} \frac{1}{\sqrt{u}}\du\]
	\noindent Now, let's apply $u$-sub,
	\[
	\text{Let: } w = \sqrt{u}, \qquad \dw = \frac{1}{2\sqrt{u}}\dw, \qquad u = w^2
	\]
	\noindent Next, substitute back in and factor out the $2$.
	\[
	\Gamma\left(\frac{1}{2}\right) = 2 \int_{w=0}^{\infty} e^{-w^2} \dw
	\]
	\noindent Looking at the graph for this new function, we can see that the integral can consume the 2 and expand the integral bounds from $0 \rightarrow \infty$ to $-\infty \rightarrow \infty$.
	\begin{figure}[!hbtp]
		\centering
		\begin{tikzpicture}
			\begin{axis}[
				axis equal,
				axis lines = center,
				xmin = -3,
				xmax = 3,
				ymin = -1.5,
				ymax = 1.5,
				minor tick num=4,
				minor tick style={draw=none},
				grid = both,
				width = \textwidth,
				height = 0.5\textwidth,
				axis line style={draw=black, thick},
				major tick length=5pt,
				ticklabel style={font=\bfseries},
				tick style={thick},
				xlabel = {$x$},
				ylabel = {$y$},
				unbounded coords=jump,
				clip=true,
				]
				\addplot[
				name path=gauss,
				red,
				line width=1.2pt,
				domain=0:4,
				samples=1600,
				] gnuplot {exp(-x**2)};
				
				\addplot[
				red,
				line width=1.2pt,
				domain=-4:0,
				samples=1600,
				] gnuplot {exp(-x**2)};
				
				\addplot[
				name path=axis,
				domain=0:4,
				] {0};
				
				\addplot[
				fill=red!30,
				fill opacity=0.6,
				] fill between[
				of=gauss and axis
				];
			\end{axis}
		\end{tikzpicture}
		\caption{Graph of the Bell Curve $f(x)=e^{-x^2}$.}
	\end{figure}
	
	
	\noindent So, let's change the bounds. Let's also set the integral to some arbitrary variable, $J$.
	\[
	\displaystyle{J = \int_{w = -\infty}^{\infty} e^{-w^2} \dw}
	\]
	\noindent Since $e^{-w^2}$ has no elementary anti derivative, we can use a Mutlivariable Calculus trick to solve the integral. First, we'll square both sides, splitting the right hand side into two integrals with arbitrary variables, $x$ and $y$.
	
	\[
	\displaystyle{J^2 = \left(\int_{x=-\infty}^{\infty}e^{-x^2} \dx\right)\left(\int_{y=-\infty}^{\infty}e^{-y^2}\dy\right)}
	\]
	\noindent Now, we need to combine them. We'll do that using Fubini's Theorem, turning the squared/multiplied integrals into a double integral.
	\[
	J^2 = \int_{x=-\infty}^{\infty}\int_{y=-\infty}^{\infty}e^{-x^2}e^{-y^2}\dy\dx
	\]
	\noindent Looking at this, the $x$ and $y$ covers the entire Cartesian Plane. So, we can compact the double integral to cover all real numbers on the 2D plane, $\mathbb{R}^2$.
	\[
	\displaystyle{
		J^2=\iint\limits_{\mathbb{R}^2}e^{-x^2}e^{-y^2}\dx\dy = \iint\limits_{\mathbb{R}^2}e^{-x^2-y^2}\dx\dy
	}
	\]
	\noindent The exponent on the $e$ looks similar to the circle equation, so we're going to convert the integral into polar coordinates.
	\[
	\displaystyle{
		r^2 = x^2 + y^2 \qquad \dx\dy = r\dr\dtheta
	}
	\]
	\noindent The angle only needs to go from $0 \rightarrow 2\pi$ to cover the entire graph, as well as $0 \rightarrow \infty$ for the radius.
	\[
	\displaystyle{
		J^2 = \int_{\theta=0}^{2\pi}\int_{r=0}^{\infty}e^{-r^2}r\dr\dtheta
	}
	\]
	\noindent Now, we can split this apart again by doing a reverse Fubini's Theorem and solving the left, simple, integral.
	\[
	\displaystyle{
	J^2 = \left(\int_{\theta=0}^{2\pi}\dtheta\right)\left(\int_{r=0}^{\infty}e^{-r^2}r\dr\right) = 2\pi\int_{r=0}^{\infty}e^{-r^2}r\dr
	}
	\]
	\noindent From here, do $u$-sub letting $v=r^2$. We're using $v$ because $u$ is already used before in the problem.
	\[
	\displaystyle{
		\text{Let: } v = r^2, \qquad \dv = 2r\dr
	}
	\]
	\noindent Substituting it in, the $2$'s and $r$'s cancel.
	\[
	\displaystyle{
		J^2 = \pi\int_{r=0}^{\infty}e^{-v}\dv
	}
	\]
	
	
	\noindent Next, solve the integral and evaluate the bounds of integration.
	\[
	\displaystyle{
		J^2 = \pi\left[\frac{e^{-v}}{-1}\right]_{r=0}^{\infty} = \pi\left[0-(-1)\right]
	}
	\]
	\noindent Finally, simplify in the brackets and take the square root of both sides, isolating J.
	\[
	\displaystyle{
		J^2 = \pi \quad \rightarrow \quad J = \sqrt{\pi}
	}
	\]
	\noindent Substituting back in the original function for J, we get:
	\[
	\displaystyle{
		\boxed{\Gamma{\left(\frac{1}{2}\right)} = \sqrt{\pi}}
	}
	\]\\ [10pt]
	
	
	\newpage
	Now, we've gone over the graph of the function as well as a relatively complex example. Let's do some problems!\\[5pt]
	
	\iftoggle{solutions}{
	First, let's convert it into terms of the Gamma Function. \\[10pt]
	}{}
	
	\iftoggle{solutions}{
		\problem{$\displaystyle{\frac{3}{2}!} = \Gamma(5/2)$} \\ [5pt]
	}{
		\problem{$\displaystyle{\frac{3}{2}!}$} \\ [5pt]
	}
	
	\iftoggle{solutions}{
		
		\noindent Now, let's split it apart and pull the $z$ out,
		\[
		\displaystyle{
			\Gamma\left(\frac{5}{2}\right) = \Gamma\left(\frac{3}{2} + 1\right) = \frac{3}{2}\Gamma\left(\frac{3}{2}\right)
		}
		\]
		\noindent So, you've lowered the value in the Gamma Function, but its not yet low enough to be something we've seen before. Let's do it again.
		\[
		\displaystyle{
			\Gamma\left(\frac{5}{2}\right) = \frac{3}{2}\Gamma\left(\frac{1}{2}+1\right) = \frac{3}{2} \cdot \frac{1}{2}\Gamma\left(\frac{1}{2}\right)
		}
		\]
		\noindent Oh, we've seen that before! $\Gamma\left(\frac{1}{2}\right) = \sqrt{\pi}$. Thus,
		\[
		\displaystyle{
			\boxed{\frac{3}{2}! = \frac{3}{2} \cdot \frac{1}{2} \cdot \sqrt{\pi} = \frac{3\sqrt{\pi}}{4}}
		}
		\]
	}{}
	
	\svs
	\ipb
	
	\iftoggle{solutions}{
		Before anything, let's expand $(x+1)^4$,\\ [10pt]
	}{}
	
	\iftoggle{solutions}{
		\problem{$\displaystyle{\int_{0}^{\infty}e^{-x}(x+1)^4\dx = \int_{0}^{\infty}e^{-x}\cdot (x^4 + 4x^3 + 6x^2 + 4x + 1) \dx}$} \\ [5pt]
	}{
		\problem{$\displaystyle{\int_{0}^{\infty}e^{-x}(x+1)^4}\dx$} \\ [5pt]
	}
	
	\iftoggle{solutions}{
		\noindent Now, distribute and split apart the integral(s). I'm also going to assign the original integral the variable R:
		\[
		\displaystyle{
		R = \int_{0}^{\infty} x^4e^{-x}\dx + 4\int_{0}^{\infty}x^3e^{-x}\dx + 6\int_{0}^{\infty}x^2e^{-x}\dx + 4\int_{0}^{\infty}xe^{-x}\dx + \int_{0}^{\infty}\dx
		}
		\]
		\noindent Remember the simple Gamma Function of a Factorial, that is:
		\[
		\displaystyle{
		n! = \int_{0}^{\infty}e^{-u}u^n \du
		}
		\]
		\noindent So, if we switch each integral into factorial form,
		\[
		\displaystyle{
		R = 4! + 4\cdot 3! + 6\cdot 2! + 4\cdot 1! + 0!
		}
		\]
		\noindent Computing this gives you your answer!
		\[
		\displaystyle{
		\boxed{R = 65}
		}
		\]
	}{}
	
	\spb
	\ipb
	
	\iftoggle{solutions}{
		\problem{$\displaystyle{ \int_{1}^{\infty} \frac{(\ln (x))^3}{x^2}\dx}$} \\ [5pt]
	}{
		\problem{$\displaystyle{ \int_{1}^{\infty} \frac{(\ln (x))^3}{x^2}\dx}$} \\ [5pt]
	}
	
	\iftoggle{solutions}{
		\noindent Let's $u$-sub first.
		\[
		\displaystyle{
			\text{Let: } u = \ln(x) \quad \Longrightarrow \quad x = e^u, \qquad \dx = e^u\du
		}
		\]
		\noindent Before substituting it back in, we need to change the bounds of integration,
		\[
		\displaystyle{
			\text{When: } x = 1 \rightarrow u = \ln(1) = 0
		}
		\]
		\[
		\displaystyle{
			\text{When: } x = \infty \rightarrow u = \ln(\infty) = \infty
		}
		\]
		\noindent We can now substitute everything in,
		\[
		\displaystyle{
			\int_{1}^{\infty} \frac{(\ln (x))^3}{x^2}\dx = \int_{0}^{\infty} \frac{u^3}{e^{2u}} \cdot e^u \du
		}
		\]
		\noindent We can now cancel things and move values up to the numerator, since we're looking for the Gamma Function form.
		\[
		\displaystyle{
			\int_{1}^{\infty} \frac{(\ln (x))^3}{x^2}\dx = \int_{0}^{\infty} u^3e^{-u} \du
		}
		\]
		\noindent That's a recognizable form
		\[
		\displaystyle{
			\boxed{\int_{1}^{\infty} \frac{(\ln (x))^3}{x^2}\dx = 3! = 6}
		}
		\]
	}{}
	
	\spb
	\ipb
	
	\iftoggle{solutions}{
		\problem{$\displaystyle{ \int_{0}^{\infty} x^6e^{-4x^2}\dx }$}
	}{
		\problem{$\displaystyle{ \int_{0}^{\infty} x^6e^{-4x^2}\dx }$}
	}
	
	\iftoggle{solutions}{
		\noindent First step, we'll perform a $u$-sub,
		\[
		\displaystyle{
			\text{Let: } u = 4x^2, \qquad \frac{\du}{8x} = \dx, \qquad x^2 = \frac{u}{4}, \qquad x = \frac{\sqrt{u}}{2}
		}
		\]
		\noindent Don't forget to also check the bounds,
		\[
		\displaystyle{
			\text{When: } x=0 \Longrightarrow u = 0, \qquad x = \infty \Longrightarrow u = \infty
		}
		\]
		\noindent We also need to split apart the $x^6$ to a defined value we already have.
		\[
		\displaystyle{
			x^6 = \left(x^2\right)^3, \qquad x^6 = \left(\frac{u}{4}\right)^3 = \frac{u^3}{64}
		}
		\]
		\noindent There we go! Let's substitute in!
		\[
		\displaystyle{
			\int_{0}^{\infty} x^6e^{-4x^2}\dx = \int_{0}^{\infty} \frac{u^3}{64} e^{-u} \frac{1}{4 \sqrt{u}} \du = \frac{1}{256}\int_{0}^{\infty}u^{\frac{5}{2}} e^{-u}\du
		}
		\]
		\noindent Recall the Gamma Function:
		\[
		\displaystyle{
			\Gamma(a) = \int_{0}^{\infty} u^{a-1}e^{-u}\du
		}
		\]
		\noindent So, if we manipulate the exponent for $u$, we can get it into the right form and condense the expression into Gamma form ($\Gamma(a)$).
		\[
		\displaystyle{
			\int_{0}^{\infty} x^6e^{-4x^2}\dx = \frac{1}{256}\int_{0}^{\infty}u^{\frac{7}{2} - 1} e^{-u}\du = \frac{1}{256} \cdot \Gamma\left(\frac{7}{2}\right)
		}
		\]
		\noindent Now, we have seen something similar with fractions, like in Problem 1. Let's use the recurrence relation to break it down.
		\[
		\displaystyle{
			\Gamma\left(\frac{7}{2}\right) = \Gamma\left(\frac{5}{2} + 1\right) = \frac{5}{2}\Gamma\left(\frac{5}{2}\right) = \dots = \frac{5}{2}\cdot\frac{3}{2}\cdot\frac{1}{2}\Gamma\left(\frac{1}{2}\right)
		}
		\]
		\noindent But, we know that $\Gamma\frac{1}{2} = \sqrt{\pi}$, so plugging everything in now,
		\[
		\displaystyle{
			\int_{0}^{\infty} x^6e^{-4x^2}\dx = \frac{1}{256}\cdot\frac{5}{2}\cdot\frac{3}{2}\cdot\frac{1}{2}\cdot\sqrt{\pi}
		}
		\]
		\noindent Finally, that gives you:
		\[
		\displaystyle{
			\boxed{\int_{0}^{\infty} x^6e^{-4x^2}\dx = \frac{15 \sqrt{\pi}}{2048}}
		}
		\]
	
	
	}{}
	
	\spb
	\ipb
	
	\iftoggle{solutions}{
		First, we'll pull the $x$ out to simplify the radical. Also use your log rules to simplify the argument of the natural log. \\ [10pt]
	}{}
	
	\iftoggle{solutions}{
		\problem{$\displaystyle{ \int_{0}^{1} x^3 \sqrt{x \ln\left(\frac{1}{x}\right)}\dx = \int_{0}^{1} x^3 \sqrt{x}\sqrt{-\ln\left(x\right)}\dx = \int_{0}^{1} x^{\frac{7}{2}} \sqrt{-\ln\left(x\right)}\dx}$} \\ [5pt]
	}{
		\problem{$\displaystyle{ \int_{0}^{1} x^3 \sqrt{x \ln\left(\frac{1}{x}\right)}\dx }$} \\ [5pt]
	}
	
	\iftoggle{solutions}{
		\noindent Let's do a $u$-sub to start,
		\[
		\displaystyle{
		\text{Let: } u = -\ln(x) \quad \Longrightarrow \quad x=e^{-u}, \qquad \dx = -e^{-u}\du
		}
		\]
		\noindent Adjust the bounds,
		\[
		\displaystyle{
			\text{When: } x=0 \Longrightarrow u = \infty, \qquad x=1 \Longrightarrow u = 0
		}
		\]
		\noindent Substitute in all the values,
		\[
		\displaystyle{
			\int_{0}^{1} x^3 \sqrt{x \ln\left(\frac{1}{x}\right)}\dx = \int_{\infty}^{0} e^{-\frac{7}{2}u} \sqrt{u} \left(-e^{-u}\right)\du
		}
		\]
		\noindent Flip the bounds and let the negatives cancel,
		\[
		\displaystyle{
			\int_{0}^{1} x^3 \sqrt{x \ln\left(\frac{1}{x}\right)}\dx = \int_{0}^{\infty} e^{-\frac{9}{2}u}u^{\frac{1}{2}}\du
		}
		\]
		\noindent We're going to do another substitution, because we need the exponent on the $e$ to be $-u$.
		\[
		\displaystyle{
			\text{Let: }t = \frac{9}{2}u \quad \Longrightarrow \quad u = \frac{2}{9}t, \qquad \dt = \frac{9}{2}\du
		}
		\]
		\noindent Now, substitute it all in, simplify the constants and factor them out.
		\[
		\displaystyle{
			\int_{0}^{1} x^3 \sqrt{x \ln\left(\frac{1}{x}\right)}\dx = \frac{2}{9}\int_{0}^{\infty}e^{-t}\left(\frac{2}{9}t\right)^{\frac{1}{2}}\dt = \frac{2\sqrt{2}}{27}\int_{0}^{\infty}e^{-t}t^{\frac{1}{2}}\dt
		}
		\]
		\noindent Recognize the Gamma Function, and use recurrence relation,
		\[
		\displaystyle{
			\int_{0}^{1} x^3 \sqrt{x \ln\left(\frac{1}{x}\right)}\dx = \frac{2\sqrt{2}}{27}\int_{0}^{\infty}e^{-t}t^{\frac{3}{2}-1}\dt = \frac{2\sqrt{2}}{27}\Gamma\left(\frac{3}{2}\right) = \frac{2\sqrt{2}}{27}\cdot\frac{1}{2}\cdot\sqrt{\pi}
		}
		\]
		\noindent Finally, cancel the $2$'s and you're finished!
		\[
		\displaystyle{
			\boxed{\int_{0}^{1} x^3 \sqrt{x \ln\left(\frac{1}{x}\right)}\dx = \frac{\sqrt{2\pi}}{27}}
		}
		\]
	}{}
\end{document}